\chapter{2008年新课标卷（理）}

\section{单选题}

\begin{problem}
已知函数 \(y = 2\sin (\omega x + \varphi)\)（\(\omega > 0\)）在区间 \([0, 2\pi]\) 的图象如图，那么 \(\omega =\)（\quad）

\bitmapfigure[max width=0.42\linewidth]{img/2008/new_curriculum_science/q1_fig1.png}

\choices
    {\(1\)}
    {\(2\)}
    {\(\frac{1}{2}\)}
    {\(\frac{1}{3}\)}
\begin{answer}
B
\end{answer}

\begin{solution}
由图象可见，在区间 \([0,2\pi]\) 内函数完成两个完整周期，所以其周期为 \(T=\pi\). 对于函数 \(y=2\sin(\omega x+\varphi)\)，周期满足 \(T=\frac{2\pi}{\omega}\)，故 \(\frac{2\pi}{\omega}=\pi\)，解得 \(\omega=2\). 因而选 B.
\end{solution}
\end{problem}

\begin{problem}
已知复数 \(z = 1 - \i\)，则 \(\frac{z^2 - 2z}{z - 1} =\)（\quad）

\choices
    {\(2\i\)}
    {\(-2\i\)}
    {\(2\)}
    {\(-2\)}
\begin{answer}
B
\end{answer}

\begin{solution}
由 \(z=1-\i\) 得 \(z-1=-\i\)，且
\(z^2=\left(1-\i\right)^2=-2\i\). 因此
\(z^2-2z=-2\i-2\left(1-\i\right)=-2\)，从而
\[
\frac{z^2-2z}{z-1}=\frac{-2}{-\i}=\frac{2}{\i}=-2\i
\]
故选 B.
\end{solution}
\end{problem}

\begin{problem}
如果等腰三角形的周长是底边长的 \(5\) 倍，那么它的顶角的余弦值为（\quad）

\choices
    {\(\frac{5}{18}\)}
    {\(\frac{3}{4}\)}
    {\(\frac{\sqrt{3}}{2}\)}
    {\(\frac{7}{8}\)}
\begin{answer}
D
\end{answer}

\begin{solution}
设等腰三角形的底边长为 \(c\)，两腰长均为 \(l\). 由周长条件，\(2l+c=5c\)，所以 \(l=2c\). 设顶角为 \(\theta\)，由余弦定理得
\[
\cos\theta=\frac{l^2+l^2-c^2}{2l^2}=\frac{2\left(2c\right)^2-c^2}{2\left(2c\right)^2}=\frac{7}{8}
\]
故选 D.
\end{solution}
\end{problem}

\begin{problem}
设等比数列 \(\{a_{n}\}\) 的公比 \(q = 2\)，前 \(n\) 项和为 \(S_{n}\)，则 \(\frac{S_4}{a_2} =\)（\quad）

\choices
    {\(2\)}
    {\(4\)}
    {\(\frac{15}{2}\)}
    {\(\frac{17}{2}\)}
\begin{answer}
C
\end{answer}

\begin{solution}
设等比数列首项为 \(a_1\). 因为公比 \(q=2\)，所以 \(a_2=2a_1\)，并且
\[
S_4=a_1\left(1+2+2^2+2^3\right)=15a_1
\]
由于题中商有意义，\(a_1\ne0\)，故
\[
\frac{S_4}{a_2}=\frac{15a_1}{2a_1}=\frac{15}{2}
\]
故选 C.
\end{solution}
\end{problem}

\begin{problem}
下面的程序框图，如果输入三个实数 \(a\)，\(b\)，\(c\)，要求输出这三个数中最大的数，那么在空白的判断框中，应该填入下面四个选项中的（\quad）

\bitmapfigure[max width=0.42\linewidth]{img/2008/new_curriculum_science/q5_fig1.png}

\choices
    {\(c > x\)}
    {\(x > c\)}
    {\(c > b\)}
    {\(b > c\)}
\begin{answer}
A
\end{answer}

\begin{solution}
程序先令 \(x=a\). 第一个判断框若填 \(b>x\)，则在 \(b>x\) 时令 \(x=b\)，否则保留 \(x=a\)，因此经过这一判断后 \(x=\max\{a,b\}\). 第二个判断框必须在 \(c>x\) 时令 \(x=c\)，否则保持原值，才能保证最后输出 \(\max\{a,b,c\}\). 所以应填 \(c>x\)，故选 A.
\end{solution}
\end{problem}

\begin{problem}
已知 \(a_{1} > a_{2} > a_{3} > 0\)，则使得 \((1 - a_{i}x)^{2} < 1\)（\(i = 1\)，\(2\)，\(3\)）都成立的 \(x\) 取值范围是（\quad）

\choices
    {\(\left(0, \frac{1}{a_1}\right)\)}
    {\(\left(0, \frac{2}{a_1}\right)\)}
    {\(\left(0, \frac{1}{a_3}\right)\)}
    {\(\left(0, \frac{2}{a_3}\right)\)}
\begin{answer}
B
\end{answer}

\begin{solution}
对任意 \(i\in\{1,2,3\}\)，有
\[
\left(1-a_i x\right)^2<1
\Longleftrightarrow -1<1-a_i x<1
\Longleftrightarrow 0<a_i x<2
\]
因为 \(a_i>0\)，所以该不等式等价于 \(0<x<\frac{2}{a_i}\). 三个不等式同时成立时，\(x\) 必须属于这些区间的交集；又因 \(a_1>a_2>a_3>0\)，有 \(\frac{2}{a_1}<\frac{2}{a_2}<\frac{2}{a_3}\)，故交集为 \(\left(0,\frac{2}{a_1}\right)\). 故选 B.
\end{solution}
\end{problem}

\begin{problem}
\(\frac{3 - \sin 70^{\circ}}{2 - \cos^{2}10^{\circ}} =\)（\quad）

\choices
    {\(\frac{1}{2}\)}
    {\(\frac{\sqrt{2}}{2}\)}
    {\(2\)}
    {\(\frac{\sqrt{3}}{2}\)}
\begin{answer}
C
\end{answer}

\begin{solution}
利用诱导公式 \(\sin70^{\circ}=\cos20^{\circ}\)，以及二倍角公式
\(\cos^2 10^{\circ}=\frac{1+\cos20^{\circ}}{2}\)，可得
\[
2-\cos^2 10^{\circ}=2-\frac{1+\cos20^{\circ}}{2}=\frac{3-\cos20^{\circ}}{2}
\]
因此
\[
\frac{3-\sin70^{\circ}}{2-\cos^2 10^{\circ}}
=\frac{3-\cos20^{\circ}}{\frac{3-\cos20^{\circ}}{2}}=2
\]
故选 C.
\end{solution}
\end{problem}

\begin{problem}
平面向量 \(\bs{a}\)，\(\bs{b}\) 共线的充要条件是（\quad）

\choices
    {\(\bs{a}\)，\(\bs{b}\) 方向相同}
    {\(\bs{a}\)，\(\bs{b}\) 两向量中至少有一个为零向量}
    {\(\exists \lambda \in \R\)，\(\bs{b} = \lambda \bs{a}\)}
    {存在不全为零的实数 \(\lambda_{1}\)，\(\lambda_{2}\)，\(\lambda_{1} \bs{a} + \lambda_{2} \bs{b} = 0\)}
\begin{answer}
D
\end{answer}

\begin{solution}
若存在不全为零的实数 \(\lambda_1\)，\(\lambda_2\)，使得
\(\lambda_1\bs a+\lambda_2\bs b=\bs0\)，则 \(\bs a\)，\(\bs b\) 线性相关；在平面向量中，这等价于两个向量共线. 反之，若 \(\bs a\)，\(\bs b\) 共线，则当其中一个为零向量时可直接取相应系数为 \(1\)、另一个为 \(0\)；当两者均非零时，一个向量是另一个向量的实数倍，也能得到这样的非零系数组合. 因此选项 D 是充要条件，故选 D.
\end{solution}
\end{problem}

\begin{problem}
甲，乙，丙\(3\text{ 位}\)志愿者安排在周一至周五的\(5\text{ 天}\)中参加某项志愿者活动. 要求每人参加一天且每天至多安排一人，并要求甲安排在另外两位前面，不同的安排方法共有（\quad）

\choices
    {\(20\text{ 种}\)}
    {\(30\text{ 种}\)}
    {\(40\text{ 种}\)}
    {\(60\text{ 种}\)}
\begin{answer}
A
\end{answer}

\begin{solution}
先从 \(5\) 天中选出甲、乙、丙参加活动的三天，共有 \(\mathrm C_5^3=10\) 种选法. 对于选定的三天，甲必须排在乙、丙之前，所以甲只能占这三天中最早的一天，乙、丙还可以互换，有 \(2\) 种排法. 因而安排方法共有
\[
\mathrm C_5^3\times2=10\times2=20
\]
种，故选 A.
\end{solution}
\end{problem}

\begin{problem}
由直线 \(x = \frac{1}{2}\)，\(x = 2\)，曲线 \(y = \frac{1}{x}\) 及 \(x\) 轴所围图形的面积为（\quad）

\choices
    {\(\frac{15}{4}\)}
    {\(\frac{17}{4}\)}
    {\(\frac{1}{6}\ln 2\)}
    {\(2 \ln 2\)}
\begin{answer}
D
\end{answer}

\begin{solution}
在区间 \(\left[\frac12,2\right]\) 上，曲线 \(y=\frac1x\) 位于 \(x\) 轴上方，所围面积为
\[
S=\int_{1/2}^{2}\frac{1}{x}\,\mathrm dx
=\left.\ln x\right|_{1/2}^{2}
=\ln2-\ln\frac12=2\ln2
\]
故选 D.
\end{solution}
\end{problem}

\begin{problem}
已知点 \(P\) 在抛物线 \(y^{2} = 4x\) 上，那么点 \(P\) 到点 \(Q(2, -1)\) 的距离与点 \(P\) 到抛物线焦点距离之和取得最小值时，点 \(P\) 的坐标为（\quad）

\choices
    {\(\left(\frac{1}{4}, -1\right)\)}
    {\(\left(\frac{1}{4}, 1\right)\)}
    {\((1,2)\)}
    {\((1,-2)\)}
\begin{answer}
A
\end{answer}

\begin{solution}
设抛物线上点 \(P\) 的参数坐标为 \(P\left(t^2,2t\right)\). 抛物线 \(y^2=4x\) 的焦点为 \(F\left(1,0\right)\)，准线为 \(x=-1\). 令准线上与 \(P\) 同纵坐标的点为 \(D\left(-1,2t\right)\)，则抛物线的定义给出 \(PF=PD=t^2+1\). 因而
\[
PQ+PF=PQ+PD\ge QD=\sqrt{9+\left(2t+1\right)^2}\ge3
\]
当且仅当 \(t=-\frac12\) 时，\(Q\)、\(P\)、\(D\) 共线且 \(QD=3\)，此时
\[
P=\left(t^2,2t\right)=\left(\frac14,-1\right)
\]
所以距离和的最小值在 \(P\left(\frac14,-1\right)\) 处取得，故选 A.
\end{solution}
\end{problem}

\begin{problem}
某几何体的一条棱长为 \(\sqrt{7}\)，在该几何体的正视图中，这条棱的投影是长为 \(\sqrt{6}\) 的线段. 在该几何体的侧视图与俯视图中，这条棱的投影分别是长为 \(a\) 和 \(b\) 的线段，则 \(a + b\) 的最大值为（\quad）

\choices
    {\(2\sqrt{2}\)}
    {\(2\sqrt{3}\)}
    {\(4\)}
    {\(2 \sqrt{5}\)}
\begin{answer}
C
\end{answer}

\begin{solution}
把这条棱在三个互相垂直方向上的分量平方分别记为 \(u\)、\(v\)、\(w\). 棱长为 \(\sqrt7\)，所以 \(u+v+w=7\). 正视图中的投影长为 \(\sqrt6\)，说明其中两个方向的分量平方之和为 \(6\)，另一个分量平方为 \(1\). 设前者为 \(u\)，\(v\)，则 \(u+v=6\)，且两个另外的正投影长度分别为
\(a=\sqrt{u+1}\)，\(b=\sqrt{v+1}\)
（交换 \(u\)，\(v\) 不影响结论）. 由均方不等式，
\[
a+b\le\sqrt{2\left(\left(u+1\right)+\left(v+1\right)\right)}
=\sqrt{16}=4
\]
当且仅当 \(u=v=3\) 时取等号，因此 \(a+b\) 的最大值为 \(4\)，故选 C.
\end{solution}
\end{problem}

\section{填空题}

\begin{problem}
已知向量 \(\bs{a} = (0, -1, 1)\)，\(\bs{b} = (4, 1, 0)\)，\(|\lambda \bs{a} + \bs{b}| = \sqrt{29}\) 且 \(\lambda > 0\)，则 \(\lambda = \)\fillinblank{}.
\begin{answer}
3
\end{answer}

\begin{solution}
由向量坐标，
\[
\lambda\bs a+\bs b=\left(4,1-\lambda,\lambda\right)
\]
由长度条件得
\[
16+\left(1-\lambda\right)^2+\lambda^2=29
\]
即 \(\lambda^2-\lambda-6=0\)，所以 \(\lambda=3\) 或 \(\lambda=-2\). 因为 \(\lambda>0\)，故 \(\lambda=3\).
\end{solution}
\end{problem}

\begin{problem}
设双曲线 \(\frac{x^2}{9} - \frac{y^2}{16} = 1\) 的右顶点为 \(A\)，右焦点为 \(F\). 过点 \(F\) 平行双曲线的一条渐近线的直线与双曲线交于点 \(B\)，则 \(\triangle AFB\) 的面积为\fillinblank{}.
\begin{answer}
\(\frac{32}{15}\)
\end{answer}

\begin{solution}
双曲线的实半轴、虚半轴分别为 \(a=3\)、\(b=4\)，所以焦半距
\(c=\sqrt{a^2+b^2}=5\)，从而 \(A=\left(3,0\right)\)、\(F=\left(5,0\right)\). 双曲线的一条渐近线为 \(y=\frac43x\)，过 \(F\) 的平行直线可写为
\[
y=\frac43\left(x-5\right)
\]
将其代入双曲线方程，得
\[
\frac{x^2}{9}-\frac{\left(\frac43\left(x-5\right)\right)^2}{16}=1
\Longleftrightarrow x^2-\left(x-5\right)^2=9
\]
故 \(x=\frac{17}{5}\)，相应地 \(y=-\frac{32}{15}\). 另一条渐近线所得交点关于 \(x\) 轴对称，三角形面积相同. 因此
\[
S_{\triangle AFB}=\frac12\cdot AF\cdot\left|y_B\right|
=\frac12\cdot2\cdot\frac{32}{15}=\frac{32}{15}
\]
\end{solution}
\end{problem}

\begin{problem}
一个六棱柱的底面是正六边形，其侧棱垂直底面. 已知该六棱柱的顶点都在同一个球面上，且该六棱柱的体积为 \(\frac{9}{8}\)，底面周长为 \(3\)，则这个球的体积为\fillinblank{}.
\begin{answer}
\(\frac{4\pi}{3}\)
\end{answer}

\begin{solution}
设正六边形边长为 \(s\)，六棱柱高为 \(h\). 由底面周长为 \(3\)，得 \(s=\frac12\). 正六边形面积为
\[
S=\frac{3\sqrt3}{2}s^2=\frac{3\sqrt3}{8}
\]
由体积条件 \(Sh=\frac98\)，得 \(h=\sqrt3\). 正六边形的外接圆半径等于边长，所以球心位于两底面中心连线的中点，球半径满足
\[
R^2=s^2+\left(\frac h2\right)^2=\frac14+\frac34=1
\]
故球的体积为
\[
V=\frac43\pi R^3=\frac{4\pi}{3}
\]
\end{solution}
\end{problem}

\begin{problem}
从甲，乙两品种的棉花中各抽测了 \(25\text{ 根}\) 棉花的纤维长度（单位：\(\mathrm{mm}\)），结果如下：

\begin{center}
\begin{tabular}{c|ccccc}
\hline
品种 & \multicolumn{5}{c}{纤维长度（单位：\(\mathrm{mm}\)）} \\
\hline
甲品种 & \(271\) & \(273\) & \(280\) & \(285\) & \(285\) \\
 & \(287\) & \(292\) & \(294\) & \(295\) & \(301\) \\
 & \(303\) & \(303\) & \(307\) & \(308\) & \(310\) \\
 & \(314\) & \(319\) & \(323\) & \(325\) & \(325\) \\
 & \(328\) & \(331\) & \(334\) & \(337\) & \(352\) \\
\hline
乙品种 & \(284\) & \(292\) & \(295\) & \(304\) & \(306\) \\
 & \(307\) & \(312\) & \(313\) & \(315\) & \(315\) \\
 & \(316\) & \(318\) & \(318\) & \(320\) & \(322\) \\
 & \(322\) & \(324\) & \(327\) & \(329\) & \(331\) \\
 & \(333\) & \(336\) & \(337\) & \(343\) & \(356\) \\
\hline
\end{tabular}
\end{center}

由以上数据设计了如下茎叶图：
\begin{center}
\begin{tabular}{rrrrr|c|lllllll}
\multicolumn{5}{c|}{甲} & & \multicolumn{7}{c}{乙} \\
\hline
 &  &  & \(3\) & \(1\) & \(27\) &  &  &  &  &  &  &  \\
 & \(7\) & \(5\) & \(5\) & \(0\) & \(28\) & \(4\) &  &  &  &  &  &  \\
 &  & \(5\) & \(4\) & \(2\) & \(29\) & \(2\) & \(5\) &  &  &  &  &  \\
\(8\) & \(7\) & \(3\) & \(3\) & \(1\) & \(30\) & \(4\) & \(6\) & \(7\) &  &  &  &  \\
 &  & \(9\) & \(4\) & \(0\) & \(31\) & \(2\) & \(3\) & \(5\) & \(5\) & \(6\) & \(8\) & \(8\) \\
 & \(8\) & \(5\) & \(5\) & \(3\) & \(32\) & \(0\) & \(2\) & \(2\) & \(4\) & \(7\) & \(9\) &  \\
 &  & \(7\) & \(4\) & \(1\) & \(33\) & \(1\) & \(3\) & \(6\) & \(7\) &  &  &  \\
 &  &  &  &  & \(34\) & \(3\) &  &  &  &  &  &  \\
 &  &  &  & \(2\) & \(35\) & \(6\) &  &  &  &  &  &  \\
\end{tabular}
\end{center}

根据以上茎叶图，对甲，乙两品种棉花的纤维长度作比较，写出两个统计结论：

\begin{circlelist}
\item\fillinblank{}；
\item\fillinblank{}.
\end{circlelist}

\begin{answer}
乙品种棉花的纤维平均长度大于甲品种棉花的纤维平均长度；甲品种棉花的纤维长度较乙品种棉花的纤维长度更分散.
\end{answer}

\begin{solution}
由原始数据求和，甲、乙两品种的纤维长度总和分别为 \(7682\) 和 \(7975\)，故样本平均数分别为
\(\bar x_{\text{甲}}=\frac{7682}{25}=307.28\)，\(\bar x_{\text{乙}}=\frac{7975}{25}=319\).
所以乙品种棉花的纤维平均长度大于甲品种. 再以极差衡量离散程度，甲、乙两品种的极差分别为
\(352-271=81\)，\(356-284=72\).
因此甲品种棉花的纤维长度比乙品种的更分散. 例如可写出两个统计结论：乙品种的平均纤维长度较大；甲品种的纤维长度较分散.
\end{solution}
\end{problem}

\section{解答题}

\begin{problem}
已知 \(\{a_{n}\}\) 是一个等差数列，且 \(a_{2} = 1\)，\(a_{5} = -5\).

\begin{enumerate}
\item 求 \(\{a_{n}\}\) 的通项 \(a_{n}\)；
\item 求 \(\{a_{n}\}\) 前 \(n\) 项和 \(S_{n}\) 的最大值.
\end{enumerate}

\begin{solution}
\begin{enumerate}
\item 设等差数列的公差为 \(d\)，则
\(a_2=a_1+d=1\)，\(a_5=a_1+4d=-5\).
两式相减得 \(3d=-6\)，所以 \(d=-2\)，再由 \(a_1+d=1\) 得 \(a_1=3\). 因而
\[
a_n=a_1+(n-1)d=3-2(n-1)=5-2n
\]
\item 由等差数列前 \(n\) 项和公式，
\[
S_n=\frac{n\left(a_1+a_n\right)}{2}=\frac{n\left(3+5-2n\right)}{2}=n(4-n)=4-(n-2)^2
\]
因为 \(n\) 是正整数，\((n-2)^2\ge0\)，且当 \(n=2\) 时取到最小值 \(0\)，所以 \(S_n\) 的最大值为 \(4\)，在 \(n=2\) 时取得.
\end{enumerate}
\end{solution}
\end{problem}

\begin{problem}
如图，已知点 \(P\) 在正方体 \(ABCD - A'B'C'D'\) 的对角线 \(BD'\) 上，\(\angle PDA = 60^\circ\).

\begin{enumerate}
\item 求 \(DP\) 与 \(CC'\) 所成角的大小；
\item 求 \(DP\) 与平面 \(A^{\prime}AD^{\prime}D\) 所成角的大小.
\end{enumerate}

\bitmapfigure[max width=0.38\linewidth]{img/2008/new_curriculum_science/q18_fig1.png}

\begin{solution}
\begin{enumerate}
\item 设正方体棱长为 \(a\)，建立直角坐标系
\(A\left(0,0,0\right)\)，\(\ B\left(a,0,0\right)\)，\(\ D\left(0,a,0\right)\)，\(\ D'\left(0,a,a\right)\).
令 \(P\) 将 \(BD'\) 按参数 \(t\) 表示为
\(P=\left(a\left(1-t\right),at,at\right)\)，\(0\le t\le1\).
于是
\(\overrightarrow{DP}=a\left(1-t,-(1-t),t\right)\)，\(\overrightarrow{DA}=a\left(0,-1,0\right)\).
由 \(\angle PDA=60^{\circ}\)，得
\[
\frac{\overrightarrow{DP}\cdot\overrightarrow{DA}}{|DP|\,|DA|}
=\frac{1-t}{\sqrt{2\left(1-t\right)^2+t^2}}=\frac12
\]
所以 \(t^2=2\left(1-t\right)^2\)，从而 \(t=\sqrt2\left(1-t\right)\). 又
\[
|DP|=a\sqrt{2\left(1-t\right)^2+t^2}=2a(1-t)
\]
因此
\[
\cos\angle\left(DP,CC'\right)=\frac{at}{|DP|}=\frac{\sqrt2}{2}
\]
故 \(DP\) 与 \(CC'\) 所成角为 \(45^{\circ}\).
\item 平面 \(A'AD'D\) 的方程为 \(x=0\)，其法向量沿 \(x\) 轴. 设 \(DP\) 与该平面所成角为 \(\alpha\)，则
\[
\sin\alpha=\frac{a(1-t)}{|DP|}=\frac12
\]
因为 \(0\le\alpha\le90^{\circ}\)，所以 \(\alpha=30^{\circ}\).
\end{enumerate}
\end{solution}
\end{problem}

\begin{problem}
\(A\)，\(B\) 两个投资项目的利润率分别为随机变量 \(X_{1}\) 和 \(X_{2}\). 根据市场分析，\(X_{1}\) 和 \(X_{2}\) 的分布列分别为

\begin{center}
\begin{tabular}{c|cc}
\hline
\(X_1\) & \(5\%\) & \(10\%\) \\
\hline
\(P\) & \(0.8\) & \(0.2\) \\
\hline
\end{tabular}
\qquad
\begin{tabular}{c|ccc}
\hline
\(X_2\) & \(2\%\) & \(8\%\) & \(12\%\) \\
\hline
\(P\) & \(0.2\) & \(0.5\) & \(0.3\) \\
\hline
\end{tabular}
\end{center}

\begin{enumerate}
\item 在 \(A\)，\(B\) 两个项目上各投资 \(100\text{ 万元}\)，\(Y_{1}\) 和 \(Y_{2}\) 分别表示投资项目 \(A\) 和 \(B\) 所获得的利润，求方差 \(D(Y_1)\)，\(D(Y_2)\)；
\item 将 \(x\text{ 万元}\)（\(0 \leqslant x \leqslant 100\)）投资 \(A\) 项目，\((100-x)\text{ 万元}\)投资 \(B\) 项目，\(f(x)\) 表示投资 \(A\) 项目所得利润的方差与投资 \(B\) 项目所得利润的方差的和. 求 \(f(x)\) 的最小值，并指出 \(x\) 为何值时，\(f(x)\) 取到最小值.（注：\(D(aX+b)=a^{2}D(X)\)）
\end{enumerate}

\begin{solution}
\begin{enumerate}
\item 将百分数按小数表示，则
\[
E(X_1)=0.05\times0.8+0.10\times0.2=0.06
\]
\[
D(X_1)=0.8\left(0.05-0.06\right)^2+0.2\left(0.10-0.06\right)^2=0.0004
\]
由于 \(Y_1=100X_1\)，所以
\[
D(Y_1)=100^2D(X_1)=4
\]
对第二个项目，
\[
E(X_2)=0.02\times0.2+0.08\times0.5+0.12\times0.3=0.08
\]
\[
D(X_2)=0.2\left(0.02-0.08\right)^2+0.5\left(0.08-0.08\right)^2+0.3\left(0.12-0.08\right)^2=0.0012
\]
由于 \(Y_2=100X_2\)，所以 \(D(Y_2)=100^2D(X_2)=12\).
\item 投资 \(A\) 项目 \(x\) 万元、投资 \(B\) 项目 \(100-x\) 万元时，两项投资所得利润的方差分别为 \(x^2D(X_1)\) 与 \((100-x)^2D(X_2)\). 因而
\[
f(x)=0.0004x^2+0.0012(100-x)^2
=\frac{x^2+3(100-x)^2}{2500}
\]
整理得
\[
f(x)=\frac{4(x-75)^2+7500}{2500}=\frac{(x-75)^2}{625}+3
\]
在 \(0\le x\le100\) 时，\((x-75)^2\ge0\)，所以 \(f(x)\) 的最小值为 \(3\)，且当 \(x=75\) 时取得.
\end{enumerate}
\end{solution}
\end{problem}

\begin{problem}
在直角坐标系 \(xOy\) 中，椭圆 \(C_1\)：\(\frac{x^2}{a^2} + \frac{y^2}{b^2} = 1\)（\(a > b > 0\)）的左，右焦点分别为 \(F_1\)，\(F_2\)，\(F_2\) 也是抛物线 \(C_2\)：\(y^2 = 4x\) 的焦点，点 \(M\) 为 \(C_1\) 与 \(C_2\) 在第一象限的交点，且 \(|MF_2| = \frac{5}{3}\).

\begin{enumerate}
\item 求 \(C_1\) 的方程；
\item 平面上的点 \(N\) 满足 \(\overrightarrow{MN} = \overrightarrow{MF_1} + \overrightarrow{MF_2}\)，直线 \(l\parallel MN\)，且与 \(C_1\) 交于 \(A\)，\(B\) 两点，若 \(\overrightarrow{OA} \cdot \overrightarrow{OB} = 0\)，求直线 \(l\) 的方程.
\end{enumerate}

\begin{solution}
\begin{enumerate}
\item 抛物线 \(C_2:y^2=4x\) 的焦点为 \(F_2\left(1,0\right)\). 设 \(M=(x_M,y_M)\)，因 \(M\) 在第一象限且在 \(C_2\) 上，有 \(y_M^2=4x_M\). 于是
\[
|MF_2|=\sqrt{\left(x_M-1\right)^2+y_M^2}=\sqrt{\left(x_M-1\right)^2+4x_M}=x_M+1
\]
由 \(|MF_2|=\frac53\)，得 \(x_M=\frac23\)，再得 \(y_M=\frac{2\sqrt6}{3}\). 椭圆的左焦点为 \(F_1=(-1,0)\)，并且
\[
|MF_1|=\sqrt{\left(x_M+1\right)^2+y_M^2}=\frac73
\]
所以椭圆的长轴长 \(2a=|MF_1|+|MF_2|=4\)，即 \(a=2\). 又焦半距 \(c=1\)，故 \(b^2=a^2-c^2=3\)，从而
\[
C_1:\frac{x^2}{4}+\frac{y^2}{3}=1
\]
\item 由焦点坐标关于原点对称，\(\overrightarrow{MF_1}+\overrightarrow{MF_2}=-2\overrightarrow{OM}\)，所以 \(\overrightarrow{MN}=-2\overrightarrow{OM}\)，从而 \(N=-M\). 因此直线 \(l\parallel MN\) 的斜率为
\[
k=\frac{y_M}{x_M}=\sqrt6
\]
可设 \(l:y=\sqrt6x+r\). 设其与椭圆交于 \(A(x_1,y_1)\)、\(B(x_2,y_2)\)，代入椭圆方程得
\[
\frac94x^2+\frac{2\sqrt6}{3}rx+\frac{r^2}{3}-1=0
\]
故由韦达定理
\(x_1+x_2=-\frac{8\sqrt6 r}{27}\)，\(x_1x_2=\frac{4r^2}{27}-\frac49\).
条件 \(\overrightarrow{OA}\cdot\overrightarrow{OB}=0\) 给出
\[
0=x_1x_2+y_1y_2
=x_1x_2+\left(\sqrt6x_1+r\right)\left(\sqrt6x_2+r\right)
=7x_1x_2+\sqrt6r(x_1+x_2)+r^2
\]
代入上式，得 \(\frac{7r^2}{27}-\frac{28}{9}=0\)，所以 \(r^2=12\)，即 \(r=\pm2\sqrt3\). 因而
\[
l:y=\sqrt6x-2\sqrt3\quad\text{或}\quad y=\sqrt6x+2\sqrt3
\]
\end{enumerate}
\end{solution}
\end{problem}

\begin{problem}
设函数 \(f(x) = ax + \frac{1}{x + b}\)（\(a\)，\(b \in \Z\)），曲线 \(y = f(x)\) 在点 \((2, f(2))\) 处的切线方程为 \(y = 3\).

\begin{enumerate}
\item 求 \(f(x)\) 的解析式；
\item 证明：函数 \(y = f(x)\) 的图象是一个中心对称图形，并求其对称中心；
\item 证明：曲线 \(y = f(x)\) 上任一点的切线与直线 \(x = 1\) 和直线 \(y = x\) 所围三角形的面积为定值，并求出此定值.
\end{enumerate}

\begin{solution}
\begin{enumerate}
\item 因切线方程为 \(y=3\)，切点为 \(\left(2,3\right)\)，故 \(f(2)=3\)，且 \(f'(2)=0\). 由
\[
f'(x)=a-\frac{1}{\left(x+b\right)^2}
\]
得 \(a=\frac{1}{\left(b+2\right)^2}\). 令整数 \(t=b+2\)，则 \(a=1/t^2\) 为整数，所以 \(t=\pm1\)，且 \(a=1\). 再由 \(f(2)=2+1/t=3\)，得 \(t=1\)，即 \(b=-1\). 因而
\(f(x)=x+\frac{1}{x-1}\)，\(x\ne1\).
\item 对任意 \(x\ne1\)，令 \(u=x-1\)，则
\[
f(x)=1+u+\frac1u
\]
当点 \(\left(x,f(x)\right)\) 关于 \(\left(1,1\right)\) 中心对称后变为 \(\left(2-x,2-f(x)\right)\)，而
\[
f(2-x)=1-u-\frac1u=2-f(x)
\]
所以对称点仍在曲线上，函数图象是中心对称图形，对称中心为 \(\left(1,1\right)\).
\item 设曲线上任一点为 \(P\left(1+u,1+u+\frac1u\right)\)，其中 \(u\ne0\). 此点处切线斜率为
\[
m=f'(1+u)=1-\frac1{u^2}
\]
切线与直线 \(x=1\) 的交点为
\[
C=\left(1,1+\frac2u\right)
\]
与直线 \(y=x\) 的交点为
\[
D=\left(2u+1,2u+1\right)
\]
两条固定直线的交点为 \(E=\left(1,1\right)\). 因此 \(EC=\frac2{|u|}\)，点 \(D\) 到直线 \(x=1\) 的距离为 \(2|u|\)，三角形 \(ECD\) 的面积为
\[
S=\frac12\cdot\frac2{|u|}\cdot2|u|=2
\]
故该三角形面积恒为 \(2\).
\end{enumerate}
\end{solution}
\end{problem}

\section{选考题：解答题}

\begin{problem}
如图，过圆 \(O\) 外一点 \(M\) 作它的一条切线. 切点为 \(A\)，过 \(A\) 点作直线 \(AP\) 垂直直线 \(OM\)，垂足为 \(P\).

\begin{enumerate}
\item 证明： \(OM \cdot OP = OA^2\)；
\item \(N\) 为线段 \(AP\) 上一点，直线 \(NB\) 垂直直线 \(ON\)，且交圆 \(O\) 于 \(B\) 点. 过 \(B\) 点的切线交直线 \(ON\) 于 \(K\). 证明： \(\angle OKM = 90^{\circ}\).
\end{enumerate}

\bitmapfigure[max width=0.48\linewidth]{img/2008/new_curriculum_science/q22_fig1.png}

\begin{solution}
\begin{enumerate}
\item 因为 \(OA\perp AM\)，且 \(AP\perp OM\)，所以三角形 \(OAM\) 与三角形 \(OPA\) 都是直角三角形. 又它们在点 \(O\) 处有相同的夹角，故
\(\triangle OAM\sim\triangle OPA\). 由对应边关系
\[
\frac{OP}{OA}=\frac{OA}{OM}
\]
得 \(OM\cdot OP=OA^2\).
\item 取 \(OM\) 为 \(x\) 轴，设圆半径为 \(r=OA\)，并记 \(O=(0,0)\)、\(M=(m,0)\)、\(P=(p,0)\). 由第（1）问，\(mp=r^2\). 因为 \(AP\perp OM\)，可设 \(N=(p,n)\). 由
\[
\overrightarrow{NB}\perp\overrightarrow{ON}
\]
得
\[
\overrightarrow{OB}\cdot\overrightarrow{ON}=|ON|^2
\]
圆在点 \(B\) 处的切线方程为 \(\overrightarrow{OB}\cdot\overrightarrow{OX}=r^2\). 由于 \(K\) 在直线 \(ON\) 上，设 \(\overrightarrow{OK}=\lambda\overrightarrow{ON}\)，代入切线方程并利用上式，得
\(\lambda|ON|^2=r^2\)，\(\overrightarrow{OK}=\frac{r^2}{|ON|^2}\overrightarrow{ON}\).
于是
\(|OK|^2=\frac{r^4}{|ON|^2}\)，\(\overrightarrow{OM}\cdot\overrightarrow{OK}=\frac{r^2mp}{|ON|^2}=\frac{r^4}{|ON|^2}=|OK|^2\).
所以 \(\overrightarrow{KO}\cdot\overrightarrow{KM}=0\)，即 \(\angle OKM=90^{\circ}\).
\end{enumerate}
\end{solution}
\end{problem}

\begin{problem}
已知曲线 \(C_{1}\)：\(\begin{cases}
x = \cos \theta,\\
y = \sin \theta.
\end{cases}\) （\(\theta\) 为参数），曲线 \(C_{2}\)：\(\begin{cases}
x = \frac{\sqrt{2}}{2} t - \sqrt{2},\\
y = \frac{\sqrt{2}}{2} t,
\end{cases}\) （\(t\) 为参数）.

\begin{enumerate}
\item 指出 \(C_{1}\)，\(C_{2}\) 各是什么曲线，并说明 \(C_{1}\) 与 \(C_{2}\) 公共点的个数；
\item 若把 \(C_{1}\)，\(C_{2}\) 上各点的纵坐标都压缩为原来的一半，分别得到曲线 \(C_{1}^{\prime}\)，\(C_{2}^{\prime}\). 写出 \(C_{1}^{\prime}\)，\(C_{2}^{\prime}\) 的参数方程. \(C_{1}^{\prime}\) 与 \(C_{2}^{\prime}\) 公共点的个数和 \(C_{1}\) 与 \(C_{2}\) 公共点的个数是否相同？说明你的理由.
\end{enumerate}

\begin{solution}
\begin{enumerate}
\item 由 \(x=\cos\theta\)、\(y=\sin\theta\)，消去参数得 \(x^2+y^2=1\)，所以 \(C_1\) 是单位圆. 对 \(C_2\)，由 \(y=\frac{\sqrt2}{2}t\) 得 \(t=\sqrt2y\)，代入 \(x\) 式得 \(x-y+\sqrt2=0\)，所以 \(C_2\) 是直线. 将 \(x=y-\sqrt2\) 代入圆方程，得
\[
2y^2-2\sqrt2y+1=0
\]
其判别式为 \(8-8=0\)，所以直线与圆相切，\(C_1\) 与 \(C_2\) 有一个公共点.
\item 纵坐标压缩为原来的一半后，直接得到
\[
C_1':\begin{cases}
x=\cos\theta,\\
y=\frac12\sin\theta,
\end{cases}
\qquad
C_2':\begin{cases}
x=\frac{\sqrt2}{2}t-\sqrt2,\\
y=\frac{\sqrt2}{4}t.
\end{cases}
\]
消去参数后，\(C_1'\) 满足 \(x^2+4y^2=1\)，\(C_2'\) 满足 \(x-2y+\sqrt2=0\). 联立得
\[
8y^2-4\sqrt2y+1=0
\]
其判别式仍为 \(32-32=0\)，故压缩后仍有一个公共点. 这也符合将平面作同一纵向压缩变换的性质：该变换是一一对应的，原两曲线的公共点恰好一一对应到压缩后两曲线的公共点，所以公共点个数相同.
\end{enumerate}
\end{solution}
\end{problem}

\begin{problem}
已知函数 \(f(x) = |x - 8| - |x - 4|\).

\begin{enumerate}
\item 作出函数 \( y = f(x) \) 的图象；
\item 解不等式 \( |x - 8| - |x - 4| > 2 \).
\end{enumerate}

\FigureLayoutDeclare{img/2008/new_curriculum_science/q24_fig1.png}{6.0cm}{2bp 2.4bp 1.2bp 0bp}
\bitmapfigure[width=6.0cm]{img/2008/new_curriculum_science/q24_fig1.png}

\begin{solution}
\begin{enumerate}
\item 按绝对值定义分三种情况：
\[
f(x)=\begin{cases}
\left(8-x\right)-\left(4-x\right)=4,&x\le4,\\
\left(8-x\right)-\left(x-4\right)=-2x+12,&4<x\le8,\\
\left(x-8\right)-\left(x-4\right)=-4,&x>8.
\end{cases}
\]
所以函数图象由直线 \(y=4\) 在 \(x\le4\) 上的部分、连接点 \(\left(4,4\right)\) 与 \(\left(8,-4\right)\) 的线段 \(y=-2x+12\)，以及直线 \(y=-4\) 在 \(x>8\) 上的部分组成. 以 \(\left(4,4\right)\)、\(\left(8,-4\right)\) 为转折点即可在题图坐标系中作出图象.
\item 同样分情况讨论. 当 \(x\le4\) 时，\(f(x)=4>2\)，所以得到 \(x\le4\). 当 \(4<x\le8\) 时，
\[
-2x+12>2\Longleftrightarrow x<5
\]
与本区间合并得 \(4<x<5\). 当 \(x>8\) 时，\(f(x)=-4\)，不满足不等式. 因而解集为
\[
\left(-\infty,4\right]\cup\left(4,5\right)=\left(-\infty,5\right)
\]
\end{enumerate}
\end{solution}
\end{problem}
